\documentclass[10pt,a4paper]{amsart}
\usepackage[margin=19mm]{geometry}
\usepackage{amsmath,amssymb,mathtools}
\usepackage[scr=rsfs,cal=euler]{mathalfa}
\usepackage{xfrac}
\usepackage[unicode,hidelinks,bookmarks=false]{hyperref}
\newcommand{\ie}{i.e.\ }
\newcommand{\cf}{cf.\ }
\newcommand{\resp}{respectively\ }
\newcommand{\Subsection}{\subsection}
\providecommand{\nobreakdash}{\nobreak\hbox{-}}
\setlength{\emergencystretch}{2em}
\multlinegap0pt
\usepackage{amssymb}
\usepackage{mathtools}
\newtheorem{lemma}{Lemma}
\title{Hopf $2$-cocycles for certain affine algebraic reductive groups}
\author{Shlomo Gelaki}
\date{Source: arXiv:2607.18599v2 (CC BY 4.0)}
\begin{document}
\maketitle

\noindent\emph{Source and licence.} Hopf $2$-cocycles for certain affine algebraic reductive groups. Version arXiv:2607.18599v2 (CC BY 4.0), distributed by arXiv under the Creative Commons Attribution 4.0 International licence, http://creativecommons.org/licenses/by/4.0/, as recorded in the arXiv licence metadata for this exact version and shown on the abstract page https://arxiv.org/abs/2607.18599v2. Full licence notice: \texttt{licenses/arxiv-2607.18599-CC-BY-4.0.txt}.\par\smallskip

\noindent\textit{Editorial scope.} Counted unit: the article's lemma lem:transitive-groupoid-reduction (source0.tex lines 263--289). The article's complete original proof of it is delivered with it (source0.tex lines 291--346, one \texttt{proof} environment of 56 lines). No mathematics is written, repaired or re-derived: the statement and the proof are the article's own lines, in the article's own notation, and no retained line is substituted or re-flowed.  No further source-local context is needed: every cross-reference of every retained line resolves inside the two delivered spans.  Nothing was omitted from any retained span, so the package carries no omission record.  No retained line contains a citation, so the wrapper carries no bibliography and no bibliography entry.  No retained line contains a figure environment, an image-inclusion command or drawing code, so no image is delivered and the package declares no asset dependency.  Editorial formatting only, in addition: an AMS \texttt{amsart} preamble that restates the article's own declarations and macros used by the retained lines, namely the environments \texttt{lemma} on separate counters, together with the standard packages those lines need.  The retained environments print in this document as Lemma 1 in order of appearance, whereas the article prints the same items with the numbering of its own full document.  The complete native source of the article as distributed in the arXiv e-print for this exact version is bundled unchanged as \texttt{sources/205662/source0.tex}.
\begin{lemma}[Reduction of a transitive action groupoid]
\label{lem:transitive-groupoid-reduction}
Assume that the action of $\Gamma$ on the finite set $X$ is transitive.  Fix
$x_0\in X$, let
\[
 H=\Gamma_{x_0}:=\{g\in\Gamma\mid gx_0=x_0\},
 \qquad
 \omega_0(g,h):=\omega(g,h;x_0),\quad g,h\in H,
\]
and set $\mathcal A:=\mathbb C^{\omega}[\Gamma\ltimes X]$.  Then
\[
 e_{x_0}\mathcal A e_{x_0}\cong\mathbb C^{\omega_0}[H].
\]
Moreover, $e_{x_0}$ is a full idempotent and, after choosing one arrow from
$x_0$ to every object of $X$, there is an algebra isomorphism
\begin{equation}\label{eq:transitive-groupoid-matrix-reduction}
 \mathcal A
 \cong
 \operatorname{Mat}_{|X|}\!\left(\mathbb C^{\omega_0}[H]\right).
\end{equation}
Consequently, restriction to the corner gives an equivalence
\[
 {\rm Mod}(\mathcal A)
 \simeq
 {\rm Mod}(\mathbb C^{\omega_0}[H]).
\]
\end{lemma}

\begin{proof}
The corner $e_{x_0}\mathcal A e_{x_0}$ is spanned by the loops $[h,x_0]$ with
$h\in H$, and their multiplication is
\[
 [g,x_0][h,x_0]=\omega_0(g,h)[gh,x_0].
\]
This proves the identification with $\mathbb C^{\omega_0}[H]$.

For each $x\in X$, choose $r_x\in\Gamma$ with $r_xx_0=x$, taking
$r_{x_0}=1$, and let
\[
 p_x:=[r_x,x_0].
\]
The two products
\[
 [r_x^{-1},x]p_x
 \quad\text{and}\quad
 p_x[r_x^{-1},x]
\]
are nonzero scalar multiples of $e_{x_0}$ and $e_x$, respectively.  The cocycle
identity shows that the two scalars are equal.  Hence one can rescale
$[r_x^{-1},x]$ to an arrow $q_x$ such that
\begin{equation}\label{eq:connecting-arrows}
 q_xp_x=e_{x_0},
 \qquad
 p_xq_x=e_x.
\end{equation}
In particular, $e_x=p_xe_{x_0}q_x$ belongs to
$\mathcal A e_{x_0}\mathcal A$ for every $x$, so $e_{x_0}$ is full.

Let $E_{x',x}$ denote the standard matrix unit indexed by $x',x\in X$.  For
$a\in e_{x'}\mathcal A e_x$, define
\begin{equation}\label{eq:explicit-groupoid-matrix-map}
 \Phi(a)=E_{x',x}\otimes q_{x'}ap_x.
\end{equation}
The middle factor belongs to $e_{x_0}\mathcal A e_{x_0}$.  Using
\eqref{eq:connecting-arrows}, one checks that
\[
 \Phi(ab)=\Phi(a)\Phi(b)
\]
whenever $a$ and $b$ are composable, while both sides vanish otherwise.  The
inverse map is
\[
 E_{x',x}\otimes b\longmapsto p_{x'}bq_x,
 \qquad
 b\in e_{x_0}\mathcal A e_{x_0}.
\]
This proves \eqref{eq:transitive-groupoid-matrix-reduction}; the stated Morita
equivalence follows, for example, from the mutually inverse functors
\[
 N\longmapsto e_{x_0}N,
 \qquad
 W\longmapsto \mathcal A e_{x_0}
       \otimes_{e_{x_0}\mathcal A e_{x_0}}W.
\]
\end{proof}

\end{document}
